\documentclass[../../../include/open-logic-section]{subfiles}

\begin{document}

\olfileid{sth}{replacement}{extrinsic}
\olsection[बाह्य विचार]{प्रतिस्थापनाविषयी बाह्य विचार}

आधी प्रतिस्थापनाचे \emph{बाह्य} समर्थन करण्याचा एक प्रयत्न पाहू. बूलोस यांनी असा युक्तिवाद सुचवला आहे:
\begin{quote}
  [\ldots] प्रतिस्थापनाची स्वयंसिद्धके स्वीकारण्याचे कारण अगदी सोपे आहे: त्यांचे अनेक इष्ट परिणाम आहेत आणि (दिसते तसे) अनिष्ट परिणाम नाहीत. क्रमिक संकल्पनेविषयीच्या प्रमेयांखेरीज त्या परिणामांत अनंत संख्यांची आदर्श नसली तरी समाधानकारक उपपत्ती आणि सुस्थापित संबंधांवरील विगमनाधारित व्याख्यांना समर्थन देणारा अत्यंत उपयुक्त निष्कर्ष येतो.
  \citep[229]{Boolos1971}
\end{quote}
बूलोस यांच्या कल्पनेचा सारांश असा की प्रतिस्थापनाचे समर्थन त्याच्या फलितांवरून करावे. त्यांनी सांगितलेली विशिष्ट फलिते आपण मागील काही प्रकरणांत पाहिली आहेत. सुक्रमणाच्या क्रमप्रकारासाठी फोन न्यूमन क्रमसूचक संख्या उत्तम प्रतिनिधी ठरतात, हे प्रतिस्थापनामुळे सिद्ध करता आले (हीच आपली ``अनंत संख्यांची आदर्श नसली तरी समाधानकारक उपपत्ती''). प्रतिस्थापनामुळे $V_\alpha$ परिभाषित करता आले, प्रतांकाची संकल्पना मांडता आली आणि $\in$-विगमन सिद्ध करता आले (हीच ``क्रमिक संकल्पनेविषयीची प्रमेये''). अखेरीस, प्रतिस्थापनामुळे अनंतपार पुनरावर्तनाचे प्रमेय सिद्ध करता येते (हीच ``सुस्थापित संबंधांवरील विगमनाधारित व्याख्या'').

हे परिणाम खरोखर इष्ट आहेत. पण इतके इष्ट परिणाम प्रतिस्थापनाला \emph{समर्थन} देण्यासाठी पुरेसे आहेत का? \emph{नाही}. निदान सरळपणे तरी नाही.

एक सोपी अडचण अशी आहे: प्रतिस्थापनाचे काही इष्ट परिणाम आपण मांडले, पण त्यांपैकी अनेक परिणाम दुसऱ्या मार्गांनीही मिळू शकले असते. ही बाब जितकी माहीत असायला हवी तितकी नाही; म्हणून परिस्थिती स्पष्ट करू.

संचांची एक सोपी उपपत्ती आहे: स्तर उपपत्ती, संक्षेपाने $\LT$.\footnote{$\LT$ ची पहिली रूपे \citet{Montague1965} आणि \citet{Scott1974} यांनी दिली. \citet{Potter2004} यांनी ती सोपी करून पुस्तकरूपात सविस्तर मांडली; अलीकडे \citet{ButtonLT1} यांनी $\LT$ आणखी सोपी केली.} $\LT$ ची स्वयंसिद्धके फक्त विस्तारात्मकता, विभक्तीकरण आणि प्रत्येक संच कोणत्या तरी \emph{स्तराचा} उपसंच असतो हे विधान इतकीच आहेत. येथे ``स्तर'' अशी कल्पक व्याख्या आहे की हे स्तर आपल्या परिचयाच्या $V_\alpha$ प्रमाणे वागतात. म्हणून $\ZF$ वरून $\LT$ सिद्ध होते; पण $\LT$ ही $\ZF$ पेक्षा \emph{खूप} दुर्बल आहे. खरे तर $\LT$ वरून जोड्या, घातसंच, अनंतता किंवा प्रतिस्थापन यांपैकी काहीही मिळत नाही. $\LT$ मध्ये अनंतता आणि घातसंच जोडून मिळणाऱ्या उपपत्तीला $\Zr$ म्हणू. यावरून जोड्याही मिळतात; म्हणून $\Zr$ निदान $\Z$ इतकी सामर्थ्यवान आहे. पण प्रत्यक्षात $\Zr$ ही $\Z$ पेक्षा काटेकोरपणे अधिक सामर्थ्यवान आहे; कारण प्रत्येक संचाला प्रतांक असतो हे विधान ती जोडते (म्हणूनच तिला $\Zr$ म्हणावे असे मी सुचवतो). $\Zr$ वरून क्रमसूचक संख्यांची पूर्णपणे समाधानकारक उपपत्ती, पदानुक्रमाचे सुक्रमित टप्प्यांत विभाजन करणारे निष्कर्ष, $\in$-विगमनाचा पुरावा आणि अनंतपार पुनरावर्तनाचे एक \emph{रूप} मिळते.

थोडक्यात, बूलोस यांना ही बाब माहीत नसली, तरी त्यांनी सांगितलेले सर्व इष्ट परिणाम प्रतिस्थापनाशिवाय \emph{मिळू शकले असते}; त्यासाठी $\Z$ ऐवजी $\Zr$ वापरणे पुरेसे होते.

(मग आपण प्रचलित पद्धत वापरून $\LT$ आणि $\Zr$ ऐवजी $\ZF$ का शिकवली? दोन कारणे आहेत. पहिले: निव्वळ ऐतिहासिक कारणांमुळे $\LT$ पासून सुरुवात करणे अपारंपरिक आहे; संच उपपत्तीवरील अधिक प्रमाणित चर्चा वाचता याव्यात अशी आपली इच्छा होती. दुसरे: $\LT$ आणि $\Zr$ समजून घेण्याची तयारी झाल्यावर \citealt{Potter2004} आणि \citealt{ButtonLT1} थेट वाचता येतील.)

अर्थात $\Zr$ ही $\ZF$ पेक्षा काटेकोरपणे दुर्बल असल्याने, $\ZF$ सिद्ध करते पण $\Zr$ अनिर्णित ठेवते असे निष्कर्ष आहेत. त्यांपैकी एखाद्या निष्कर्षाच्या आधारे प्रतिस्थापनाचे बाह्य समर्थन करण्याचा प्रयत्न करता येईल. पण $\Zr$ कशी वापरावी हे समजल्यावर (अ) प्रतिस्थापनामुळेच, दुसऱ्या मार्गाने नाही, ठरणाऱ्या आणि (ब) सहज अर्थाने खऱ्या वाटणाऱ्या विधानांची अनेक उदाहरणे शोधणे बरेच कठीण असते. (अधिक माहितीसाठी \citealt[\S13.2]{Potter2004} पाहा.)

मुख्य निष्कर्ष असा: प्रतिस्थापनाचे पटणारे बाह्य समर्थन द्यायचे असेल, तर प्रतिस्थापनाशिवाय \emph{मिळूच शकणार नाही} असा निष्कर्ष शोधावा लागेल. हे सोपे काम नाही.

प्रतिस्थापनाच्या निव्वळ बाह्य समर्थनाच्या कोणत्याही प्रयत्नासमोर आणखी एक अडचण उभी राहते; कदाचित ती पहिल्यापेक्षा अधिक मूलभूत आहे. बूलोस फक्त प्रतिस्थापनाचे अनेक इष्ट परिणाम सांगत नाहीत; त्याचे ``(दिसते तसे) अनिष्ट'' परिणाम नाहीत, असेही म्हणतात. पण कंसातील ``दिसते तसे'' ही सावधगिरीची नोंद अतिशय महत्त्वाची आहे.

आपण येथे कसे पोहोचलो ते आठवू. अनौपचारिक गुणधर्माधारित संचरचनेत विसंगती निर्माण झाली; त्याला प्रतिसाद म्हणून आपण संचांची संचयी-क्रमिक संकल्पना स्वीकारली. या संकल्पनेबरोबर एक मांडणी येते आणि ती आपल्याला तिच्या सुसंगततेची खात्री देईल, अशी आशा आहे. पण त्या मांडणीच्या आतून प्रतिस्थापनाचे समर्थन करता आले नाही, तर $\ZF$ सुसंगत आहे असे मानण्याचे (अजून) कोणतेही कारण आपल्याकडे नाही. किंवा अधिक नेमके: आत्तापर्यंत कोणालाही \emph{विसंगती सापडलेली नाही} या (कदाचित निव्वळ योगायोगाच्या) वस्तुस्थितीखेरीज $\ZF$ च्या सुसंगततेवर विश्वास ठेवण्याचे कारण नाही. नेमक्या याच अर्थाने बूलोस यांचे विधान ``(दिसते तसे) $\ZF$ सुसंगत आहे'' इतकेच ठरते. सुसंगततेबाबत याहून अधिक खात्री आपण मागितली पाहिजे.

$\ZF$ च्या (ज्ञात) परिणामांपुरतेच मर्यादित असलेले, म्हणजे प्रतिस्थापनाचे \emph{निव्वळ} बाह्य समर्थन करण्याचे, कोणतेही प्रयत्न या अडचणीला सामोरे जातील. म्हणून आपण प्रतिस्थापनाचे \emph{आंतरिक} समर्थन शोधू: संचांविषयी आपण सांगितलेली मांडणी आपल्याला आधीपासूनच प्रतिस्थापन स्वीकारण्यास बांधील करते, असे सुचवणारे समर्थन.

\end{document}
